\chapter{放射监测与核查计划}

\section{辐射源监测}

建设单位已开展过后装治疗项目，本项目拟依托现有相关辐射源监测计划（如表~\ref{tab:ch5-1-t1-76007769}），并利旧部分并计划新增配备自主监测设备和相应辐射源自主监测记录表格（如表~\ref{tab:ch5-1-t2-fb61305b} 所示）。



\begin{table}[H]
\centering
\caption{本项目辐射源监测计划及评价}
\label{tab:ch5-1-t1-76007769}
\begin{tabular}{|>{\Centering}m{1.2cm}|>{\Centering}m{2.4cm}|>{\Centering}m{2.6cm}|>{\Centering}m{2.2cm}|>{\Centering}m{4.2cm}|>{\Centering}m{1cm}|}
\hline
\multirow{2}{=}{监测\\类型}
& \multicolumn{3}{c|}{计划情况}
& \multirow{2}{*}{法规/标准要求}
& \multirow{2}{*}{评价} \\
\cline{2-4}
& 监测项目 & 监测频度 & 监测方式 & & \\
\hline
验收检测
& \multirow{2}{=}{\Centering\arraybackslash 按照国家有关标准规定执行}
& 设备新安装、维修或更换重要部件后进行
& 委托有资质的单位
& 本项目应当进行职业病危害控制效果评价
& 符合 \\
\cline{1-1} \cline{3-6}
状态检测
& \hspace{0pt}
& 使用后：1次/年
& 委托有资质的单位
& 由省级卫生健康行政部门资质认证的检测机构每年至少进行一次状态检测
& 符合 \\
\hline
稳定性检测
& 按国家相关要求参数进行
& 相关参数监测周期要求
& 自主监测
& 对放射诊疗设备定期进行稳定性检测
& 符合 \\
\hline
\end{tabular}
\end{table}



\begin{table}[H]
\centering
\caption{自主监测设备配备计划及评价}
\label{tab:ch5-1-t2-fb61305b}
\begin{tabular}{|>{\Centering}m{2.4cm}|>{\Centering}m{2.6cm}|>{\Centering}m{3cm}|>{\Centering}m{0.8cm}|>{\Centering}m{3.4cm}|>{\Centering}m{1cm}|}
\hline
项目 & 设备名称 & 生产厂家/型号 & 数量 & 检定/校准计划 & 评价 \\
\hline
\multirow{7}{=}{后装治疗机}
& 放疗剂量仪 & IBA DOSE1 & 1套 & 计划委托有资质的单位进行 & \multirow{8}{*}{符合} \\
\cline{2-5}
& 井型电离室 & HDR 1000PLUS & 1个 & 计划委托有资质的单位进行 & \\
\cline{2-5}
& 智能化X、γ辐射仪 & SB-1 & 1个 & / & \\
\cline{2-5}
& 个人剂量报警仪 & JB2000 & 1套 & / & \\
\cline{2-5}
& 空盒气压表 & 待定 & 1套 & / & \\
\cline{2-5}
& 出源长度测定标尺 & 待定 & 1套 & / & \\
\cline{1-5}
放射工作场所防护检测
& 辐射剂量测量仪 & JB4000型（与直线加速器机房共用） & 1个 & 计划委托有资质的单位进行 & \\
\hline
\end{tabular}
\end{table}



\begin{table}[H]
\centering
\caption{后装治疗机稳定性检测要求}
\label{tab:ch5-1-t3-4b3d1a7c}
\begin{tabular}{|m{1cm}|m{6cm}|m{3cm}|m{4cm}|}
\hline
序号 & 检测项目 & 技术要求 & 周期 \\
\hline
1 & 源活度 & ±5\% & 换源后或维修后(192Ir) \\
\hline
2 & 源传输到位精确度 & ±1mm & 换源后或维修后(192Ir) \\
\hline
3 & 贮源器表面(5cm、100cm)泄漏辐射所致周围剂量当量率 & 5cm(50μSv/h)\newline 100cm(5μSv/h) & 换源后或维修后(192Ir) \\
\hline
\end{tabular}
\end{table}


\section{放射工作场所监测}

建设单位已开展过放射治疗项目，本项目拟依托现有的放射工作场所监测设备，并依托现有《放射治疗质量制度》《辐射治疗装置维护、维修制度》等管理制度，

监测情况具体如下：

$\textcircled{1}$对新建、改建、扩建的放射工作场所进行预评价、控制效果评价，并按要求向卫生健康行政部门申请卫生审查和卫生验收。

$\textcircled{2}$按要求请有资质的放射卫生技术服务机构每年对我院所有射线装置进行一次仪器性能和放射工作场所的检测，发现问题及时整改，确保患者和工作人员的安全。

$\textcircled{3}$为确保放射治疗和诊断质量，建设单位按标准要求定期对诊疗设备性能和防护进行检测。

$\textcircled{4}$对于新安装、维修或更换重要部件后设备，按要求请有资质的放射卫生技术服务机构进行验收检测。

$\textcircled{5}$自主检测计划：物理师每月对所有装有射线装置的机房进行辐射检测，检测设备采用便携式辐射周围剂量当量检测仪，检测的关注点应符合国家相关要求。如测量结果超出国家标准，应立即停机整改。如测量结果与以往测量基准值偏大时，需查明原因，并进行相关内容完善。

综上所述，本项目放射工作场所监测计划能满足后续工作开展需要，符合《放射治疗放射防护要求》(GBZ 121-2020)。

\section{个人剂量监测}

建设单位现有放射工作人员已委托有资质的机构进行定期个人剂量监测，本项目均依托原有放射工作人员，建设单位拟对本项目所有放射工作人员按照国家要求进行个人剂量监测。


\begin{table}[H]
\centering
\caption{个人剂量监测计划及评价}
\label{tab:ch5-3-t1-0033a307}
\begin{tabular}{|>{\Centering}m{1.6cm}|>{\Centering}m{1.6cm}|>{\Centering}m{3.2cm}|>{\Centering}m{6cm}|>{\Centering}m{1cm}|}
\hline
监测项目
& \multicolumn{2}{c|}{计划情况}
& 法规/标准要求
& 评价 \\
\hline
\multirow{4}{=}{个人剂量\\监测}
& 监测人数 & 本项目所有放射工作人员
& \multirow{2}{=}{放射工作单位应当安排本单位的放射工作人员接受个人剂量监测}
& \multirow{2}{*}{符合} \\
\cline{2-3}
& 监测种类 & X、$\gamma$射线 & & \\
\cline{2-4} \cline{5-5}
& 监测周期 & 1次/3 个月
& 常规监测周期一般为 1 个月，最长不得超过 3 个月
& 符合 \\
\cline{2-4} \cline{5-5}
& 监测单位 & 委托有资质的单位（江西辐射剂量检测院有限公司）
& 个人剂量监测工作应当由具有资质的个人剂量监测技术服务机构承担
& 符合 \\
\hline
\end{tabular}
\end{table}


\section{小结}

建设单位已开展放射治疗工作多年，已制定符合要求的辐射源监测计划和监测

表格，放射工作场所监测和个人剂量监测以委托检测和自主监测相结合的形式按照国家相关要求进行监测。

建议：落实质控设备的配置，放疗剂量仪、井型电离室、辐射剂量仪、空盒气压表、长杆水银温度计等设备在使用前送至有资质的机构进行检定或校准。建设单位应及时组织相关工作人员参加培训，物理人员熟悉掌握质控自主监测设备的使用，并在日常监测过程中做好相关记录。
